% Generated draft: requires mathematical and source review.
\documentclass[11pt]{article}
\usepackage[margin=25mm]{geometry}
\usepackage{fontspec}
\setmainfont{Latin Modern Roman}
\usepackage{amsmath,amssymb,mathtools,unicode-math}
\setmathfont{Latin Modern Math}
\usepackage{enumitem,needspace,xurl,hyperref}
\setlength{\parindent}{0pt}
\setlength{\parskip}{5pt}
\setlength{\emergencystretch}{3em}
\newcommand{\N}{\mathbb{N}}
\newcommand{\Z}{\mathbb{Z}}
\newcommand{\Q}{\mathbb{Q}}
\newcommand{\R}{\mathbb{R}}
\newcommand{\C}{\mathbb{C}}
\newcommand{\F}{\mathbb{F}}
\newcommand{\A}{\mathbb{A}}
\newcommand{\PP}{\mathbb{P}}
\newcommand{\E}{\mathbb{E}}
\newcommand{\eps}{\varepsilon}
\newcommand{\dd}{\,\mathrm{d}}
\begin{document}
\section*{UnsolvedMath ID: 30001816}
\textbf{Source problem code:} catalogue entry
\subsection{Definitions and mathematical statement}
Let \(p\) be a prime, let \(F\) be a \(p\)-adic field, and let \(E/F\) be an unramified quadratic extension. Let \(n\) denote the rank parameter in the intended arithmetic fundamental lemma (AFL) setting. Let \(\gamma\) be a regular semisimple element of the unitary symmetric space, and suppose that \(\gamma\) matches an element \(g\) of the corresponding unitary group under the matching relation of that setting.

The supplied record uses \(1_S\) for the test function on the symmetric-space side, \(O'_\gamma(1_S)\) for the differentiated orbital integral, and \(\omega(\gamma)\) for the transfer factor. It uses \(N_{n-1}\) for a unitary Rapoport--Zink space and \(\Delta(N_{n-1})\) for its indicated diagonal cycle in the intended ambient space. The notation \((1,g)\Delta(N_{n-1})\) denotes the translate of that cycle by the indicated action, and \(\langle\ ,\ \rangle\) denotes the arithmetic intersection number of the two cycles.

\textbf{Source-level conjectural statement.} In the intended AFL setting, for every regular semisimple matching pair \((\gamma,g)\), does the identity
\[
O'_\gamma(1_S)
=-\omega(\gamma)
\left\langle
\Delta(N_{n-1}),
(1,g)\Delta(N_{n-1})
\right\rangle
\]
hold? The supplied target is the unrestricted AFL identity, including cases beyond the minuscule case.

\textbf{Unresolved qualifications.} The record does not specify the admissible values of \(n\), restrictions on \(p\) or \(F\), the precise symmetric space and unitary group, the matching relation, or the moduli definitions and ambient space for the cycles. It also does not define the test function, the orbital-integral derivative and its normalization, the transfer factor, or the intersection convention. In particular, any residue-field logarithm factor and its placement are unspecified. The displayed formula therefore preserves the supplied assertion; it is not a fully normalized, self-contained formulation. The geometric hypotheses ensuring that the arithmetic intersection number is defined and finite must also be recovered from the cited sources.
\subsection{Short English statement}
The arithmetic fundamental lemma asks whether a differentiated orbital integral equals a signed arithmetic intersection number of matching cycles on unitary Rapoport--Zink spaces. This record targets the identity beyond minuscule cases, but leaves essential hypotheses and normalizations unspecified. Its partially solved status is supplied dated curation and has not been independently verified here.
\subsection{Sources}
\begin{itemize}
\item Supplied local UnsolvedMath registry record: permanent ID \(30001816\), exact source problem code \texttt{catalogue entry}, title \emph{The Arithmetic Fundamental Lemma Beyond Minuscule Cases}, in the Oberwolfach Reports Open Problems collection. The record gives the proposed year as 2011 and supplies no named proposer.
\item \emph{Algebraische Zahlentheorie} (2011), Oberwolfach Reports, pp.~1709--1768, DOI \texttt{10.4171/owr/2011/31}. Supplied report link: \url{https://doi.org/10.4171/owr/2011/31}. The supplied background describes Oberwolfach Reports as documenting mathematical workshops and their open-problem discussions.
\item Supplied literature reference labeled ``Rapoport--Terstiege--Zhang AFL'': \url{https://arxiv.org/abs/1203.5827}.
\item Supplied literature reference labeled ``original Oberwolfach report'': \url{https://doi.org/10.4171/owr/2011/31}.
\item Literature assessment dated exactly \texttt{2026-08-21}: this is supplied dated curation, not an independently verified assessment. It labels the problem partially solved, reports proofs for \(n=2\) and \(n=3\) and restricted minuscule progress in the report, attributes the general minuscule case to Rapoport--Terstiege--Zhang, and describes the unrestricted AFL conjecture as unresolved. The supplied background calls this a dated literature triage derived from OWR curation and requests independent verification of time-sensitive status and cited claims.
\item The registry labels the statement \texttt{corrected\_verified} and says that its mathematical content was checked against the official Oberwolfach report, with notation, context, and question form normalized without changing the source-level claim. This is a supplied verification assertion; no source was independently visited or checked for this draft.
\end{itemize}
\end{document}
